% \section{Existing Approaches}

% In this section, we review existing approaches for nucleus core decomposition and discuss their limitations. 
% %
% We focus on the classical peeling-based framework for nucleus decomposition~\cite{firstNucleus,fastHierarchy,sotaPrveHierarchy,sotaHierarchy, nuclearTWEB}. 
% Although local update-based methods~\cite{LocalNuclear} have also been proposed to enable parallel computation, they share the same fundamental limitation of explicit $s$-clique enumeration and are therefore not discussed in detail here due to page limit.

% \begin{algorithm}[t]
\small
\caption{Peeling Framework}
\label{alg:peel_concept}
\KwIn{Graph $G$, clique sizes $0<r<s$}
\KwOut{$\kappa_s(R)$ for every $r$-clique $R$}

%-------------------------
% Step 1 : initial counts
%-------------------------
% $mark$ every $r$-clique $R \in G$ as \texttt{active};
Enumerate all $r$-cliques in $G$;

Enumerate all $s$-cliques in $G$;

$mark$ every $s$-clique $S \in G$ as \texttt{active};


$\delta(R)\leftarrow$ the number of $s$-cliques that contain $R$

$coreness \leftarrow 0$

%-------------------------
% Step 2 : peeling loop
%-------------------------
\While{there exists \texttt{active} $r$-cliques}{
    find $c=\min\{\delta(R)\mid R\text{ is \texttt{active}}\}$

    $coreness \leftarrow max(c, coreness)$

    let $\mathcal{R}=\{R\mid \delta(R)=c,\;R\text{ \texttt{active}}\}$

    \ForEach{$R\in\mathcal{R}$}{
        $\kappa_s(R)\leftarrow coreness$

        mark $R$ as \texttt{inactive};

        \ForEach{\texttt{active} $s$-clique $S$ with $R\subseteq S$}{
            mark $S$ as \texttt{inactive};

            \ForEach{\texttt{active} $r$-clique $R'\subseteq S$}{
                $\delta(R')\leftarrow\delta(R')-1$;
            }
        }
    }
}
\Return{$\kappa_s(R)$ for every $r$-clique $R \in G$}
\end{algorithm}


% \begin{algorithm}[t]
\small
% \caption{\textsc{Local $h$-Index Framework}}
% \label{alg:hindex_concept}
% \KwIn{Graph $G$, clique sizes $0<r<s$}
% \KwOut{$\\\kappa_s(R)$ for every $r$-clique $R$}

% $\\\kappa_s(R) \leftarrow |nbr_s(R)|$ for every $r$-clique $R \in G$

% $update \leftarrow$ \texttt{true}

% \While{$update$}{
%     $update \leftarrow \texttt{false}$\;

%     \ForEach{$r$-clique $R$ in $G$}{
%         $L \leftarrow [\,]$

%         \ForEach{$s$-clique $S$ with $R\subseteq S$}{
%             $m_S \leftarrow \min_{R'\subseteq S,\;R'\neq R}\kappa(R')$\;
%             append $m_S$ to $L$\;
%         }

%         $\kappa^{new}(R) \leftarrow \mathcal{H}(L)$\;
%         \If{$\kappa^{new}(R) \neq \kappa(R)$}{
%             $\kappa(R) \leftarrow \kappa^{new}(R)$\;
%             $update \leftarrow \texttt{true}$\;
%         }   
%     }
% } 

% \Return{$\\\kappa_s(R)$ for every $r$-clique $R \in G$}

% \end{algorithm}


% % Peeling是最被广泛使用的Nuclear Core Decomposition方法, 其基本思想是通过逐层剥离（peeling）图中的拥有最少s-clique的r-clique来计算每个r-clique的coreness.在每一轮中，算法都会找到所有的活跃r-clique，并将其剥离，直至没有活跃的r-clique为止。
% Algorithm~\ref{alg:peel_concept} illustrates the basic paradigm of peeling-based algorithms. 
% It computes the coreness of each $r$-clique by iteratively peeling off the $r$-cliques with the smallest number of associated $s$-cliques in each round. 
% In every iteration, the algorithm identifies all active $r$-cliques and removes them until no active $r$-cliques remain.
% % 算法\ref{alg:peel_concept}展示了Peeling的基本框架.

% The peeling framework comprises two phases: \emph{initial counting} and the \emph{peeling loop}. 
% In the initial counting phase, all $r$-cliques and $s$-cliques are enumerated and marked as active, and the number of active $s$-cliques containing each active $r$-clique is computed. 
% In the peeling loop, the algorithm iteratively identifies and removes active $r$-cliques with the smallest number of active $s$-cliques; 
% the $s$-cliques containing these $r$-cliques are then marked as inactive. 
% This process continues until no active $r$-cliques remain. 
% The coreness of each $r$-clique is assigned to the current peeling level, which corresponds to the minimum number of active $s$-cliques containing that $r$-clique.

% % O(|K_r|+|K_s|)
% \stitle{Complexity.} 
% The time complexity of the peeling framework is $O(|K_r| + |K_s|)$ \cite{sotaPrveHierarchy}, where $|K_r|$ and $|K_s|$ represent the numbers of $r$-cliques and $s$-cliques in the graph, respectively .

% \stitle{Clique Enumeration in Existing Approaches.}
% Existing methods have time and space complexities that rely on the total number of $s$-cliques in the graph. 
% However, the number of $s$-cliques grows combinatorially with $s$, leading to prohibitive computational costs. 
% As a result, current algorithms fail to scale to larger $s$ values, with state-of-the-art methods handling only small configurations (e.g., $(4,5)$-nucleus decomposition) on small or sparse datasets. 
% Hence, it is essential to revisit the algorithm design to remove this dependency and avoid explicit enumeration of $s$-cliques.
\myaddRC{
Computing these core numbers requires tracking, for each $r$-clique $R$, its \emph{active} $s$-clique support $\deg_{(s)}(R)$ (Definition~\ref{def:r-s-degree}) under a sequence of deletions.
A widely used approach in prior work is the peeling-based nucleus decomposition~\cite{firstNucleus,fastHierarchy,sotaPrveHierarchy,sotaHierarchy,nuclearTWEB}:
starting from the full graph, it iteratively removes the active $r$-cliques with the smallest current support and updates the supports of the remaining $r$-cliques.
We summarize this baseline next to fix terminology and to expose the scalability bottleneck that motivates our index-based design.

\begin{algorithm}[t]
\small
\caption{Peeling Framework}
\label{alg:peel_concept}
\KwIn{Graph $G$, clique sizes $0<r<s$}
\KwOut{$\kappa_s(R)$ for every $r$-clique $R$}

%-------------------------
% Step 1 : initial counts
%-------------------------
% $mark$ every $r$-clique $R \in G$ as \texttt{active};
Enumerate all $r$-cliques in $G$;

Enumerate all $s$-cliques in $G$;

$mark$ every $s$-clique $S \in G$ as \texttt{active};


$\delta(R)\leftarrow$ the number of $s$-cliques that contain $R$

$coreness \leftarrow 0$

%-------------------------
% Step 2 : peeling loop
%-------------------------
\While{there exists \texttt{active} $r$-cliques}{
    find $c=\min\{\delta(R)\mid R\text{ is \texttt{active}}\}$

    $coreness \leftarrow max(c, coreness)$

    let $\mathcal{R}=\{R\mid \delta(R)=c,\;R\text{ \texttt{active}}\}$

    \ForEach{$R\in\mathcal{R}$}{
        $\kappa_s(R)\leftarrow coreness$

        mark $R$ as \texttt{inactive};

        \ForEach{\texttt{active} $s$-clique $S$ with $R\subseteq S$}{
            mark $S$ as \texttt{inactive};

            \ForEach{\texttt{active} $r$-clique $R'\subseteq S$}{
                $\delta(R')\leftarrow\delta(R')-1$;
            }
        }
    }
}
\Return{$\kappa_s(R)$ for every $r$-clique $R \in G$}
\end{algorithm}


% \begin{algorithm}[t]
\small
% \caption{\textsc{Local $h$-Index Framework}}
% \label{alg:hindex_concept}
% \KwIn{Graph $G$, clique sizes $0<r<s$}
% \KwOut{$\\\kappa_s(R)$ for every $r$-clique $R$}

% $\\\kappa_s(R) \leftarrow |nbr_s(R)|$ for every $r$-clique $R \in G$

% $update \leftarrow$ \texttt{true}

% \While{$update$}{
%     $update \leftarrow \texttt{false}$\;

%     \ForEach{$r$-clique $R$ in $G$}{
%         $L \leftarrow [\,]$

%         \ForEach{$s$-clique $S$ with $R\subseteq S$}{
%             $m_S \leftarrow \min_{R'\subseteq S,\;R'\neq R}\kappa(R')$\;
%             append $m_S$ to $L$\;
%         }

%         $\kappa^{new}(R) \leftarrow \mathcal{H}(L)$\;
%         \If{$\kappa^{new}(R) \neq \kappa(R)$}{
%             $\kappa(R) \leftarrow \kappa^{new}(R)$\;
%             $update \leftarrow \texttt{true}$\;
%         }   
%     }
% } 

% \Return{$\\\kappa_s(R)$ for every $r$-clique $R \in G$}

% \end{algorithm}

\stitle{Baseline: Peeling-Based Nucleus Decomposition.}\label{sec:existing_approaches}
% 
Algorithm~\ref{alg:peel_concept} illustrates the peeling paradigm.
In the \emph{initial counting} phase, all $r$-cliques and $s$-cliques are enumerated and marked active, and each active $r$-clique $R$ is assigned its initial support $\deg_{(s)}(R)$.
In the peeling loop, the algorithm repeatedly removes the active $r$-cliques with the smallest support;
then every active $s$-clique that contains a removed $r$-clique is marked inactive, which decreases the support of the remaining $r$-cliques that it contains.
The coreness $\kappa_s(R)$ of each $r$-clique is the peeling level at which it is removed.

This framework runs in $O(|K_r| + |K_s|)$ time~\cite{sotaPrveHierarchy}, where $|K_r|$ and $|K_s|$ denote the numbers of $r$-cliques and $s$-cliques, respectively.
However, both time and space critically depend on \emph{explicitly enumerating} $s$-cliques (during initial counting and during support updates).
Since $|K_s|$ can grow combinatorially with $s$, this dependence quickly becomes prohibitive as $s$ increases, which explains why prior methods typically target only small configurations (e.g., $(4,5)$-nucleus) on small or sparse graphs.
Local update-based variants have been proposed to improve parallelism~\cite{LocalNuclear}, but they still rely on $s$-clique enumeration and thus inherit the same bottleneck.
These limitations motivate algorithmic designs that avoid explicit $s$-clique enumeration.

\stitle{Baseline Primitive: Eager Support Maintenance.}
The reference methods in our experiments differ in implementation details.
%
However, they share the same exact peeling skeleton.
%
They repeatedly remove minimum-support active $r$-cliques.
%
They then maintain exact supports for the surviving $r$-cliques.
%
The key difference lies in the \emph{support-maintenance primitive}.
%
In prior methods, support maintenance is eager.
%
When a peeled unit invalidates incident $s$-cliques, the algorithm either enumerates the affected $s$-cliques or reconstructs the affected local structure.
%
It then recomputes the supports.
%
This primitive is \emph{reconstructive}.
%
It materializes or revisits local clique structure after each update.

\stitle{From Reconstructive to Differential Maintenance.}
Our journal-level goal is not to replace the exact peeling skeleton.
%
Our goal is to replace the reconstructive maintenance primitive.
%
The \cpi gives a compact counting representation.
%
DPCM shows how to update exact supports by \emph{differential path contribution updates}.
%
That is, we update only the local path states incident to removed units.
%
We do not re-materialize the affected $s$-clique space.
%
This point is important.
%
Our novelty is not a new decomposition paradigm.
%
Our novelty is a new exact support-maintenance theory under the same decomposition paradigm.
}
